\section{OpenAI vs Anthropic：战略是组织能力的表达}

上一章说明 Coding 为什么重要，本章看公司战略。广密把 OpenAI 和 Anthropic 的分化解释为“不同组织能力的表达”。两家公司同宗同源，但 2025 年 Q1 的押注明显不同：OpenAI 更强地押注 o 系列 reasoning model、ChatGPT 消费入口和平台野心；Anthropic 更聚焦 base model、Coding、Agentic workflow 和企业市场。

\begin{figure}[H]
\centering
\includegraphics[width=0.92\textwidth]{openai-anthropic-strategy.png}
\caption{OpenAI vs Anthropic：战略差异是组织能力、价值观和产品路径的表达。自制概念图，依据 00:19:55--00:30:18 对谈内容整理。}
\end{figure}

\begin{knowledgebox}{读图：公司战略不是 PPT，而是谁有能力做什么}
左边 OpenAI 强调模型、消费入口和 o 系列；右边 Anthropic 强调安全、企业、Coding 和 Agentic 工作流。广密的分析不是简单排名，而是在问：组织里的核心人才、产品压力和领导层关注点，会怎样改变技术路线。
\end{knowledgebox}

\subsection{o 系列、消费入口和 Pre-training 的资源竞争}

本节拆解 OpenAI 的内部张力。o 系列在数学和代码 benchmark 上进步很快，很容易获得资源和组织注意力；ChatGPT 用户增长和消费互联网化也会消耗管理层精力。与此同时，Pre-training 团队人才流动和组织调整可能削弱底座训练节奏。广密担心的是，OpenAI 过早走向消费互联网公司，会让“智能上限”主线被流量和产品增长拉偏。

\begin{warningbox}{课堂提示：流量很重要，但不等于智能主线}
如果一个模型公司过早被消费入口牵引，它可能在用户增长、内容、品牌和生态上取得优势，但也可能降低对底座训练、长期研究和能力上限的专注。这个判断需要持续观察，而不能从单季度产品热度下结论。
\end{warningbox}

\subsection{Anthropic 的 Coding 机会}

本节看 Anthropic。广密对 Anthropic 的信心来自两个方面：一是它可能仍然在 Pre-training base model 上有强组织能力；二是 Claude/Sonnet 在 Coding 和 Agentic workflow 中已经形成开发者投票。Cursor 等产品大量调用 Claude 系列模型，说明 Coding 不是边缘能力，而是能产生真实 token 消耗和商业付费的场景。

\begin{importantbox}{实践经验：开发者投票比口号更硬}
如果开发者长期愿意为某个模型的 Coding 能力付费，而且工具产品把它设为默认模型，这说明模型能力已经进入真实生产力链条。Coding 场景的价值，不只在 benchmark，而在真实代码、真实报错和真实交付。
\end{importantbox}

\subsection{硅谷认知分歧：智能重要还是流量重要}

本节把公司分化上升为认知分歧。一个阵营认为智能上限仍然是一切的根，应该继续押注 base model、Pre-training、Coding、Agent 和 AI for Science；另一个阵营更看重消费入口、用户增长、默认入口和平台生态。实际世界不会非此即彼，但不同公司会因为组织、资本和产品压力在两边摆动。

\noindent\begin{tabularx}{\textwidth}{p{0.20\textwidth}p{0.34\textwidth}X}
\toprule
路线 & 关注指标 & 风险 \\
\midrule
智能优先 & base model、Coding、Agent、科学能力、长期上限 & 短期产品和收入可能慢。 \\
流量优先 & 用户规模、入口、消费产品、默认选择 & 可能稀释研究注意力。 \\
平台优先 & API、工具协议、生态、企业工作流 & 需要同时维持能力和开发者信任。 \\
\bottomrule
\end{tabularx}

\subsection{本章小结}

OpenAI 与 Anthropic 的分化说明，AI 公司竞争不是单纯模型参数竞争，而是组织能力、路线信仰、产品入口、客户结构和商业压力的综合表达。后面的 AGI roadmap 必须在这个背景下阅读。

